\clearpage
\thispagestyle{empty}
\pdfbookmark[0]{Über achromatische Doppelobjektive. Ausgabe und Quellen}{doublet-edition}
\begin{center}\Large\textbf{Über achromatische Doppelobjektive}\end{center}

Carl Friedrich Gauß’ Abhandlung über achromatische Doppelobjektive und die Aufhebung der Farbzerstreuung ist eine KI-Übertragung des deutschen Textes der \textit{Werke}, Band V (Göttingen, 1867), S.~504--508, in modernes Deutsch. Der Satz wurde mit KI erstellt. Der vollständige Text enthält die historische Anmerkung in eckigen Klammern und beide bearbeitbaren Tabellen.

Die Abhandlung erschien im November--Dezember-Heft 1817 der von B.~von Lindenau und J.~G.~F. Bohnenberger herausgegebenen \textit{Zeitschrift für Astronomie und verwandte Wissenschaften}, Bd.~4, S.~345--351. Die Zeitschriftenfassung und die Werkausgabe sind separat verlinkt. Gauß verweist auf S.~279 und 392 von Bohnenbergers Abhandlung aus dem Jahr 1816 in Bd.~1.

\medskip\noindent
\href{https://gdz.sub.uni-goettingen.de/id/PPN236006339}{Deutsche Werkausgabe (1867)}\par
\noindent\href{https://www.digitale-sammlungen.de/view/bsb10061006?page=345}{Zeitschriftenfassung (1817)}\par
\noindent\href{https://www.digitale-sammlungen.de/view/bsb10061006?page=390}{Druckfehlerverzeichnis der Zeitschrift}\par
\noindent\href{https://www.digitale-sammlungen.de/view/bsb10061003?page=287}{Bohnenberger, S.~279}\quad
\href{https://www.digitale-sammlungen.de/view/bsb10061003?page=400}{S.~392}\par
\noindent\href{https://doi.org/10.1007/978-3-642-49319-5_24}{Deutscher Nachdruck (1877)}
\clearpage
